% Written By Enoch Yu
% Downloaded from archive.enochyu.com

\documentclass{article}

\usepackage{amsmath}
\usepackage{amssymb}
\usepackage{amsthm}
\usepackage{mathtools}

\usepackage[margin=1in, headsep=15pt]{geometry}
\usepackage{lastpage}
\usepackage{fancyhdr}
\fancyhead[R]{Enoch Yu}
\fancyfoot[C]{Page \thepage \hspace{1pt} of \pageref{LastPage}}
\pagestyle{fancy}
\usepackage{parskip}

\usepackage{enumitem}
\usepackage[dvipsnames]{xcolor}
\usepackage[normalem]{ulem}
\usepackage{url}
\usepackage{hyperref}
\hypersetup{
    colorlinks=true,
    linkcolor=black,
    filecolor=magenta,      
    urlcolor=cyan,
    pdftitle={7.11.2025 AMC 10, 12 Free Class},
    pdfpagemode=FullScreen,
}
\usepackage{tabularx}
\usepackage{arydshln}
\usepackage{float}
\usepackage{pdfpages}

\usepackage{tikz}
\usetikzlibrary{calc, angles, shapes.geometric}
\usepackage{tkz-euclide}
\usepackage{pgfplots}
\pgfplotsset{compat=1.9}
\usepackage[font=small,labelfont=bf]{caption}

\newtheorem{theorem}{Theorem}[section]
\newtheorem{lemma}[theorem]{Lemma}
\newtheorem*{lemma*}{Lemma}
\newtheorem{sublemma}{Lemma}[section]
\newtheorem{proposition}{Proposition}
\newtheorem{corollary}{Corollary}[theorem]
\newtheorem{example}{Example}[section]
\newtheorem*{example*}{Example}
\newenvironment{solution}{\begin{trivlist}\item[]{\bf Solution}}{\qed \end{trivlist}}


\title{7.11.2025 AMC 10, 12 Free Class}
\author{Enoch Yu}
\date{11 July 2025}

\begin{document}

\maketitle

\subsection*{Problem 1}
Find all pairs of integers $(x, y)$ such that
$x^3 + y^3 = (x + y)^2$.

\begin{solution}
\[
(x + y)(x^2 - xy + y^2) = (x + y)^2
\]
\[
\begin{cases}
    x + y = 0 \implies x = -y \\
    x + y \neq 0 \implies x^2 - xy + y^2 = x + y
\end{cases}
\]
Since the first case provided the solution,
the second case could be observed.
\begin{align*}
    x^2 - xy + y^2 &= x + y \\
    x^2 - x(y + 1) + y^2 - y &= 0 \\
    x &= \frac{y + 1 \pm \sqrt{y^2 + 2y + 1 - 4y^2 + 4y}}{2} \\
    D
    &= -3y^2 + 6y + 1 \geq 0 \\
    &\Rightarrow 3y^2 - 6y - 1 \leq 0 \\
    -1 < \frac{3 - 2\sqrt{3}}{3} \leq
    &y \leq \frac{3 + 2\sqrt{3}}{3} < 3 \\
    \therefore
    y = 0, 1, 2 &\Rightarrow (1, 0), (2, 1), (1, 2), (2, 2)
\end{align*}
Thus, $\boxed{(x, -x), (1, 0), (2, 1), (1, 2), (2, 2)}$
are the roots.
\end{solution}

\subsection*{Problem 2}
Solve the equation $x^3 - 3x = \sqrt{x + 2}$
for all real values of $x$.

\begin{solution}
Notice that when $x = 2$, the equation satisfies.
Therefore, intuitively, the cases for $x$ could be set.
\[
\begin{cases}
    x \leq 2 \implies \sqrt{x + 2} \geq x \\
    x > 2 \implies \sqrt{x + 2} < x
\end{cases}
\]
Observe the second equation.
\begin{align*}
    x^3 - 3x - \sqrt{x + 2} &= 0 \\
    x^3 - 3x - x &< 0 \\
    x(x^2 - 4) &< 0 \\
    \therefore x < -2 \text{ or } 0 &< x < 2
\end{align*}
Due to a contradiction, the fact that the roots of the
equation are never greater than 2 could be deduced.
Moreover, since $x$ is a real number, $-2 \leq x \leq 2$.
Let $x = 2\cos\theta$. Therefore, the following equations
could be deduced.
\begin{align*}
    8\cos^3\theta - 6\theta &= \sqrt{2\cos\theta + 2} \\
    \cos3\theta &= \sqrt{\frac{\cos\theta + 1}{2}} \\
    \cos3\theta &= \cos\frac{\theta}{2} \\
    \therefore
    x &= \boxed{2, 2\cos\frac{4\pi}{5}, 2\cos\frac{4\pi}{7}}
\end{align*}
\end{solution}
\end{document}

